% %%%%%%%%%%%%%%%%%%%%%%%%%%%%%%%%%%%%%%%%%%%%%%%%%%%%%%%%%%%%%%%%%%%%%
% Section M
% %%%%%%%%%%%%%%%%%%%%%%%%%%%%%%%%%%%%%%%%%%%%%%%%%%%%%%%%%%%%%%%%%%%%%

\section{A Section}
\label{sec:section_m}

\subsection{Example of Text Edition}

Some text. \textbf{Some text.} \textit{Some text.} \emph{Some text.} \algo{Some text.} 

Unbreakable~space. Unbreakable~space. Unbreakable~space. Unbreakable~space. Unbreakable~space.Unbreakable~space. \mbox{Unbreakable box.} \mbox{Unbreakable box.} \mbox{Unbreakable box.} \mbox{Unbreakable box.} \mbox{Unbreakable box.} Be careful with \algo{$\backslash$mbox\{\}}. If you put too much text in it, \LaTeX will have difficulties to space the text correctly and generate a \emph{Warning}, as in this example.

\bigskip

Space before new paragraph.

\if\Language\EN Example of english quotation marks: ``a quote''.
\else Example of french quotation marks: \og a quote \fg{}. \fi

\subsection{Example of References}

In this template, we use the \emph{natbib} package for the bibliography (see in \textit{Configuration.tex}). You can change the bibliography style in \textit{MAIN.tex} (\algo{$\backslash$bibliographystyle\{\}}). In this template, we use a specific one imported from ACM journal templates (\textit{ACM-Reference-Format.bst}), but there is some predefined style such as \textit{acm}, \textit{ieeetr}, and \textit{plain}, the default one.

Some text with reference~\cite{Zhang2021_Gamma}. Some text. Some text with multiple references~\cite{Srikanth2017_SuperStrider, Barrett1994_CSR_BCSR_DIA_ELL, Smith2017_FROSTT, Beamer2017_GAP, JESD235B, Brandt2003, Deveci2018}.

References to publications: Paper~\ref{publ:SoA} and Paper~\ref{publ:SpDC}.